% 05_gravitational_waves.tex
% Ch.5: GW propagation + GW detection as impedance antenna

\chapter{Gravitational Waves}
\label{ch:gravitational_waves}

\section{GW Propagation as Impedance Modulation}
\label{sec:gw_propagation}

Gravitational waves are transverse inductive shear waves in the LC
lattice --- the same medium governed by the same operators.  The
Schwarzschild metric maps directly to a radial impedance profile:
\begin{align}
    \varepsilon_{eff}(r) &= \frac{\varepsilon_0}{1 - r_s/r}
    \quad\longrightarrow\quad \varepsilon \to \infty \text{ at horizon}
    \label{eq:eps_schwarzschild} \\
    \mu_{eff}(r) &= \mu_0 \cdot (1 - r_s/r)
    \quad\longrightarrow\quad \mu \to 0 \text{ at horizon}
    \label{eq:mu_schwarzschild}
\end{align}
where $r_s = 2GM/c^2$.  The gravitational impedance is therefore:
\begin{equation}
    Z(r) = \sqrt{\mu_{eff}/\varepsilon_{eff}}
         = Z_0 \cdot (1 - r_s/r)
    \label{eq:Z_grav}
\end{equation}
and the reflection coefficient at the horizon is:
\begin{equation}
    \Gamma(r \to r_s) = \frac{Z(r) - Z_0}{Z(r) + Z_0} \to -1
    \quad\text{(total reflection)}
\end{equation}
This is the SAME $\Gamma \to -1$ that governs Pauli exclusion
(particle confinement), plasma cutoff ($\omega < \omega_p$),
and the Meissner effect ($B < B_c$).

\begin{figure}[H]
    \centering
    \includegraphics[width=0.95\textwidth]{gw_schwarzschild_impedance.pdf}
    \caption{Schwarzschild impedance profile for a 30$\,M_\odot$ black hole,
    computed from the physics engine module \texttt{gw\_propagation.py}.
    \textbf{Top left:} Permittivity diverges ($\varepsilon \to \infty$)
    and permeability vanishes ($\mu \to 0$) at the horizon.
    \textbf{Top right:} Impedance $Z(r)/Z_0$ drops to zero.
    \textbf{Bottom left:} Reflection coefficient $\Gamma \to -1$
    (total reflection --- the event horizon IS the impedance boundary).
    \textbf{Bottom right:} Refractive index $n > 1$ near the mass
    (gravitational lensing = dielectric refraction).}
    \label{fig:gw_schwarzschild}
\end{figure}

\subsection{Lossless Propagation}

LIGO gravitational waves have strain
$h \sim 10^{-21}$.  The equivalent voltage per lattice cell is
$V_{GW} = h \cdot c \cdot \ell_{\text{node}} \cdot 2\pi f
\sim 10^{-19}\;\text{V}$, which is $\sim 10^{-24}$ times smaller than
$V_{\text{snap}} \approx 511\;\text{kV}$.  Far below the saturation
threshold, the lattice acts as a perfect lossless linear transmission
line.  Gravitational waves propagate to cosmological distances without
dissipation --- matching LIGO observations exactly, with zero free
parameters.

\subsection{Black Hole Echoes}

Because the event horizon is
$\Gamma = -1$ (not absorption), post-merger GW energy reflected
between the photon sphere and the horizon produces
\textit{echoes} with delay $\Delta t \approx r_s \ln(r_{ph}/(r_{ph}
- r_s)) / c$.  For a 30$\,M_\odot$ remnant, this predicts $\Delta t
\sim 0.01$--$0.1\;\text{s}$, consistent with the tentative detections
reported by Abedi et al.\ (2017).

\section{GW Detection: The Impedance Antenna}
\label{sec:gw_detection}

A gravitational wave detector is an impedance antenna.  LIGO's 4~km
Fabry--P\'erot arms form resonant cavities in the LC vacuum, where
each light bounce amplifies the GW-induced impedance modulation.
The passing GW strain $h$ perturbs the local vacuum impedance:
\begin{equation}
    \delta Z = Z_0 \cdot h
\end{equation}
The accumulated phase shift after $N$ bounces is:
\begin{equation}
    \Delta\phi = \frac{2\pi f_{GW}}{c} \cdot L \cdot N \cdot h
\end{equation}
For LIGO ($L=4$~km, $N=280$, $h=10^{-21}$, $f=100$~Hz),
$\Delta\phi \approx 2.3 \times 10^{-21}$~rad --- resolved via
homodyne readout against 750~kW circulating laser power.

The strain sensitivity is bounded by two quantum noise sources:
shot noise (phase) and radiation pressure (amplitude), whose
geometric mean is the Standard Quantum Limit:
\begin{equation}
    h_{SQL}(f) = \sqrt{h_{shot}^2 + h_{RP}^2}
\end{equation}

\begin{figure}[H]
    \centering
    \includegraphics[width=0.95\textwidth]{gw_detector_sensitivity.pdf}
    \caption{\textbf{Left:} Strain sensitivity curves for LIGO (4~km),
    LISA (2.5~Gm), and Einstein Telescope (10~km), modeled as
    impedance antennas. GW150914 ($h\sim10^{-21}$) is marked.
    \textbf{Right:} LIGO noise budget decomposition: shot noise
    dominates at high frequency, radiation pressure at low frequency,
    with the SQL minimum determining peak sensitivity.}
    \label{fig:gw_detector}
\end{figure}

\noindent
The lattice voltage ratio for LIGO GW is:
\begin{equation}
    \frac{V_{GW}}{V_{\text{snap}}} \approx 1.4 \times 10^{-28}
\end{equation}
Twenty-eight orders of magnitude below saturation.  The vacuum is
a \textit{perfect} lossless transmission line for gravitational waves,
exactly as observed.
